\documentclass[twocolumn, a4paper, 9pt]{jarticle}

\makeatletter
\def\section{\@startsection{section}{1}{\z@}{2ex plus .2ex minus .2ex}%
  {.5ex plus .2ex minus .2ex}{\large\bfseries}}
\def\thesection{\arabic{section}.}
\def\subsection{\@startsection{subsection}{1}{\z@}{.7ex plus .2ex minus .2ex}%
  {.5ex plus .2ex minus .2ex}{\normalsize\bfseries}}
\def\thesubsection{\arabic{section}.\arabic{subsection}}
\def\thefootnote{\fnsymbol{footnote}}
\makeatother

% ipsj-kansai
\setlength{\topmargin}{-10mm} % 15mm - 1in
\setlength{\headheight}{5mm}
\setlength{\headsep}{5mm}
\setlength{\oddsidemargin}{-7mm} % 18mm - 1in
\setlength{\evensidemargin}{-7mm} % 18mm - 1in
\setlength{\textheight}{247mm} % 297 - 25(top) - 25(bottom)
\setlength{\textwidth}{174mm} % 210 - 18*2
\setlength{\columnsep}{7mm}

\usepackage{fancyhdr}
\pagestyle{empty}
\thispagestyle{fancy}
\rhead[]{\fontsize{9pt}{9pt} 研究の夢を語る会@SocSEL}
\cfoot[]{}
\renewcommand{\headrulewidth}{0pt}

\begin{document}
\twocolumn[
  \begin{center}
    {
      {\bf \Large タイトル \\} % 14.4pt
      \vspace{0.5ex}
    }
    \vspace{1.5ex}
    {
      \fontsize{10.5pt}{10pt}
      \begin{tabular}{c}
        著者$^\dag$ \\
        松田 和輝
      \end{tabular}
    }
    \vspace{1.5ex}
  \end{center}
]


%%%%%%%%%%%%%%%%%%%%
\section{はじめに}
%%%%%%%%%%%%%%%%%%%%
ソフトウェア開発において、ソフトウェアの開発者自身が開発するプログラムに加えて、ライブラリと呼ばれる他の開発者が作成したプログラムを利用することがある。
ソフトウェアの開発とライブラリの開発はそれぞれ独立して行われており、ソフトウェアの開発中にライブラリのプログラムに加えられた新機能や不具合修正をソフトウェアにも適用させたい場合、その都度新しくなったライブラリのプログラムをソフトウェアに取り込む必要がある。

ある製品に変更が加えられていく各段階をバージョン、バージョンを一意に識別する表記をバージョン名と言う。
変更が加えられ新しくなったバージョンのライブラリをソフトウェアに取り込む際は、ソフトウェアが依存するライブラリに現在付与されているバージョン名より、一つ以上新しいバージョン名を付与されたものに差し替えることによってバージョン更新を行う。
この時、古いバージョンのライブラリの仕様を前提とした現在のソフトウェアと、変更が加えられ新しいバージョンとなったライブラリの仕様に食い違いがある場合、バージョン更新によってソフトウェアが正しく動作しなくなることがある。
ソフトウェア開発者はこの問題に対して、古いバージョンのライブラリを新しいバージョンと差し替えても変わらずソフトウェアが動作するか、すなわち新しいバージョンのライブラリに後方互換性があるかを判断した上でバージョン更新を行う必要がある。

ソフトウェア開発者があるバージョンのライブラリの後方互換性の有無を判断するための方法の一つとして、セマンティックバージョニング(SemVer)がある。
SemVerとは、ライブラリ開発者がバージョン名の付与を後方互換性の有無を元に規則的に行うことで、ライブラリを使用するソフトウェア開発者がバージョン名からそのライブラリの後方互換性の有無を知ることが出来るようにすることなどを目的とした、ソフトウェア開発におけるいくつかの規則のことである。
しかしながら、一般に任意の二つのプログラムの互換性の有無を正確に判断することは困難であり、SemVerの適用はライブラリ開発者が手動で行うことから、SemVerが適用されたバージョン名が正確に後方互換性の有無を表しているとは限らないという問題がある。

本研究では、単体テストに加えられる変更に着目して互換性の有無の判断することで、自動的にSemVerに基づいたバージョン名を付与する手法を提案する。

%%%%%%%%%%%%%%%%%%%%
\section{手法}
%%%%%%%%%%%%%%%%%%%%

%%%%%%%%%%%%%%%%%%%%
\section{結果（期待する結果）}
%%%%%%%%%%%%%%%%%%%%

%%%%%%%%%%%%%%%%%%%%
\section{今後の課題}
%%%%%%%%%%%%%%%%%%%%

\bibliographystyle{ipsjunsrt}
\bibliography{bibfile}

\end{document}